\documentclass[xcolor={table}]{beamer}

% --- TEMA E CORES (AZUL CLÍNICA) ---
\usetheme{Madrid}
\definecolor{ClinicBlue}{RGB}{13,71,161}
\setbeamercolor{structure}{fg=ClinicBlue}
\setbeamercolor{palette primary}{bg=ClinicBlue,fg=white}
\setbeamercolor{palette secondary}{bg=ClinicBlue!80!black,fg=white}
\setbeamercolor{palette tertiary}{bg=ClinicBlue!60!black,fg=white}
\setbeamercolor{title in head/foot}{bg=ClinicBlue,fg=white}

\usepackage[brazil]{babel}
\usepackage[utf8]{inputenc}
\usepackage[T1]{fontenc}
\usepackage{ragged2e}
\usepackage{xcolor}
\usepackage{amssymb}
\usepackage{booktabs}
\usepackage{array}
\usepackage{tikz}

\title{Clínica Asklepion}
\subtitle{Sistema Web de Gestão de Consultas Médicas}
\author{Caio Mendes}
\institute{Clínica Asklepion}
\date{2026}

\begin{document}

% ============================================================
% SLIDE 1: CAPA
% ============================================================
\begin{frame}[plain]
    \centering
    \vspace{0.8cm}
    {\color{ClinicBlue}\Huge\textbf{+}}\\
    \vspace{0.4cm}
    \begin{beamercolorbox}[sep=8pt,center,shadow=true,rounded=true]{title}
        \usebeamerfont{title}\inserttitle\par
        \vspace{2pt}
        \usebeamerfont{subtitle}\insertsubtitle\par
    \end{beamercolorbox}
    \vspace{0.5cm}
    \begin{beamercolorbox}[sep=5pt,center,rounded=true]{palette primary}
        \insertauthor\ $\cdot$\ \insertdate
    \end{beamercolorbox}
\end{frame}

% ============================================================
% SLIDE 2: SUMÁRIO
% ============================================================
\begin{frame}{Sumário}
    \begin{columns}
        \begin{column}{0.5\textwidth}
            \begin{enumerate}
                \item Objetivos
                \item Problemas
                \item Tecnologias Utilizadas
                \item Requisitos do Sistema
                \item Arquitetura do Sistema
            \end{enumerate}
        \end{column}
        \begin{column}{0.5\textwidth}
            \begin{enumerate}
                \setcounter{enumi}{5}
                \item Controle de Acesso
                \item Modelo Entidade-Relacionamento
                \item Regras de Negócio
                \item Resultados Obtidos
                \item Conclusão
            \end{enumerate}
        \end{column}
    \end{columns}
\end{frame}

% ============================================================
% SLIDE 3: OBJETIVOS
% ============================================================
\begin{frame}{Objetivos}
    \begin{block}{Objetivo Geral}
        \justifying
        Desenvolver um sistema web completo para gerenciamento de consultas médicas,
        integrando pacientes, médicos e equipe administrativa em uma plataforma única e
        acessível.
    \end{block}
    \vspace{0.3cm}
    \textbf{Objetivos Específicos:}
    \begin{itemize}
        \item[\color{ClinicBlue}$\checkmark$] Automatizar o agendamento de consultas por parte dos pacientes
        \item[\color{ClinicBlue}$\checkmark$] Prover painel dedicado para cada perfil de usuário
        \item[\color{ClinicBlue}$\checkmark$] Gerenciar disponibilidade de horários dos médicos dinamicamente
        \item[\color{ClinicBlue}$\checkmark$] Permitir cancelamento e remarcação de consultas
        \item[\color{ClinicBlue}$\checkmark$] Centralizar o controle administrativo de médicos e consultas
    \end{itemize}
\end{frame}

% ============================================================
% SLIDE 4: PROBLEMAS
% ============================================================
\begin{frame}{Problemas Identificados}
    \begin{columns}
        \begin{column}{0.55\textwidth}
            \begin{alertblock}{Processo Manual}
                Agendamentos feitos por telefone ou presencialmente, sujeitos a erros e conflitos de horários.
            \end{alertblock}
            \vspace{0.2cm}
            \begin{alertblock}{Falta de Visibilidade}
                Pacientes não tinham acesso aos horários disponíveis de cada médico em tempo real.
            \end{alertblock}
            \vspace{0.2cm}
            \begin{alertblock}{Controle Descentralizado}
                Equipe administrativa sem ferramenta unificada para acompanhar e gerir consultas.
            \end{alertblock}
        \end{column}
        \begin{column}{0.45\textwidth}
            \begin{exampleblock}{Solução Proposta}
                \justifying
                \small
                Um sistema web com \textbf{3 painéis distintos} (paciente, médico, equipe) e uma
                \textbf{API REST} centralizada que elimina redundâncias e conflitos,
                garantindo consistência dos dados em tempo real.
            \end{exampleblock}
        \end{column}
    \end{columns}
\end{frame}

% ============================================================
% SLIDE 5: TECNOLOGIAS
% ============================================================
\begin{frame}{Tecnologias Utilizadas}
    \begin{columns}
        \begin{column}{0.5\textwidth}
            \textbf{Frontend}
            \begin{itemize}
                \item \textbf{HTML5} — Estrutura das páginas
                \item \textbf{CSS3} — Estilização responsiva
                \item \textbf{JavaScript} (Vanilla) — Lógica e interatividade
                \item \texttt{localStorage} — Sessão do paciente
            \end{itemize}
            \vspace{0.3cm}
            \textbf{Infraestrutura}
            \begin{itemize}
                \item \textbf{Apache/XAMPP} — Servidor de arquivos estáticos
                \item \textbf{Node.js} — Execução do servidor
            \end{itemize}
        \end{column}
        \begin{column}{0.5\textwidth}
            \textbf{Backend}
            \begin{itemize}
                \item \textbf{Node.js} — Ambiente de execução
                \item \textbf{Express.js} — Framework HTTP/REST
                \item \textbf{SQLite3} — Banco de dados relacional
                \item \textbf{CORS + body-parser} — Middlewares
            \end{itemize}
            \vspace{0.3cm}
            \begin{block}{Endpoints da API}
                \small
                $\approx$\textbf{15 endpoints REST} cobrindo autenticação,
                médicos, consultas e disponibilidade.
            \end{block}
        \end{column}
    \end{columns}
\end{frame}

% ============================================================
% SLIDE 6: REQUISITOS FUNCIONAIS
% ============================================================
\begin{frame}{Requisitos Funcionais}
    \small
    \begin{table}
        \rowcolors{1}{white}{ClinicBlue!10}
        \begin{tabular}{>{\bfseries}p{1.1cm} p{8.2cm}}
            \toprule
            \textbf{Código} & \textbf{Descrição} \\
            \midrule
            RF01 & Autenticação por CPF e senha (paciente via local; equipe via API) \\
            RF02 & Cadastro de pacientes com validação de CPF e senha \\
            RF03 & Agendamento de consultas com seleção de médico, dia e horário \\
            RF04 & Cancelamento e remarcação de consultas pelo paciente \\
            RF05 & Painel médico: visualizar, concluir e cancelar consultas \\
            RF06 & Gerenciamento de disponibilidade de horários pelo médico \\
            RF07 & Painel admin: cadastrar, editar e remover médicos \\
            RF08 & Painel admin/equipe: concluir e cancelar qualquer consulta \\
            \bottomrule
        \end{tabular}
    \end{table}
\end{frame}

% ============================================================
% SLIDE 7: REQUISITOS NÃO FUNCIONAIS
% ============================================================
\begin{frame}{Requisitos Não Funcionais}
    \begin{columns}
        \begin{column}{0.5\textwidth}
            \begin{block}{Desempenho}
                \begin{itemize}
                    \item[\color{ClinicBlue}$\checkmark$] Slots calculados dinamicamente para os próximos 14 dias
                    \item[\color{ClinicBlue}$\checkmark$] Respostas da API em tempo real sem recarregamento de página
                \end{itemize}
            \end{block}
            \vspace{0.2cm}
            \begin{block}{Usabilidade}
                \begin{itemize}
                    \item[\color{ClinicBlue}$\checkmark$] Interface responsiva (mobile e desktop)
                    \item[\color{ClinicBlue}$\checkmark$] Navegação por abas sem redirecionamento de página
                \end{itemize}
            \end{block}
        \end{column}
        \begin{column}{0.5\textwidth}
            \begin{block}{Segurança}
                \begin{itemize}
                    \item[\color{ClinicBlue}$\checkmark$] Autenticação obrigatória para todas as áreas internas
                    \item[\color{ClinicBlue}$\checkmark$] Isolamento de dados: cada perfil acessa apenas o que lhe pertence
                \end{itemize}
            \end{block}
            \vspace{0.2cm}
            \begin{block}{Manutenibilidade}
                \begin{itemize}
                    \item[\color{ClinicBlue}$\checkmark$] Banco inicializado via \texttt{schema.sql}
                    \item[\color{ClinicBlue}$\checkmark$] Scripts de inicialização automatizados (\texttt{.bat})
                \end{itemize}
            \end{block}
        \end{column}
    \end{columns}
\end{frame}

% ============================================================
% SLIDE 8: ARQUITETURA
% ============================================================
\begin{frame}{Arquitetura do Sistema}
    \justifying
    O sistema segue a arquitetura \textbf{Cliente-Servidor em 3 camadas}:
    \vspace{0.3cm}

    \begin{center}
    \begin{tikzpicture}[
        box/.style={draw, rounded corners, fill=ClinicBlue!15, minimum width=3.2cm,
                    minimum height=0.8cm, font=\small\bfseries, text centered},
        arr/.style={->, thick, color=ClinicBlue}
    ]
        \node[box] (front) at (0,0)    {Camada de\\Apresentação};
        \node[box] (api)   at (4.5,0)  {Camada de\\Negócio (API)};
        \node[box] (db)    at (9,0)    {Camada de\\Dados (SQLite)};

        \draw[arr] (front.east) -- node[above]{\tiny HTTP/REST} (api.west);
        \draw[arr] (api.east)   -- node[above]{\tiny SQL}       (db.west);
        \draw[arr] (api.west)   -- node[below]{\tiny JSON}      (front.east);
        \draw[arr] (db.west)    -- node[below]{\tiny Resultados}(api.east);
    \end{tikzpicture}
    \end{center}

    \vspace{0.2cm}
    \begin{columns}
        \begin{column}{0.33\textwidth}
            \centering\small
            \textbf{HTML/CSS/JS}\\
            Páginas SPA-like,\\3 painéis distintos
        \end{column}
        \begin{column}{0.33\textwidth}
            \centering\small
            \textbf{Node.js/Express}\\
            API REST\\porta 3000
        \end{column}
        \begin{column}{0.33\textwidth}
            \centering\small
            \textbf{SQLite3}\\
            4 tabelas,\\arquivo local
        \end{column}
    \end{columns}
\end{frame}

% ============================================================
% SLIDE 9: CONTROLE DE ACESSO
% ============================================================
\begin{frame}{Controle de Acesso}
    \small
    \begin{table}
        \rowcolors{1}{white}{ClinicBlue!10}
        \begin{tabular}{>{\bfseries}p{2cm} p{3cm} p{4.5cm}}
            \toprule
            \textbf{Perfil} & \textbf{Autenticação} & \textbf{Permissões} \\
            \midrule
            Paciente   & CPF + senha\newline(localStorage)  & Agendar, ver, cancelar e remarcar suas consultas \\
            Médico     & CPF + senha\newline(API/banco)      & Ver suas consultas, concluir, cancelar, gerenciar disponibilidade \\
            Admin      & CPF + senha\newline(API/banco)      & Tudo do Admin: cadastrar/editar/remover médicos, gerir todas as consultas \\
            Recepção   & CPF + senha\newline(API/banco)      & Visualizar e gerir consultas (sem acesso à aba de médicos) \\
            \bottomrule
        \end{tabular}
    \end{table}
    \vspace{0.2cm}
    \begin{exampleblock}{Isolamento garantido}
        \small
        A função \texttt{requireAuth(roles)} em \texttt{auth.js} verifica o perfil na sessão e redireciona automaticamente caso o acesso não seja permitido.
    \end{exampleblock}
\end{frame}

% ============================================================
% SLIDE 10: MER
% ============================================================
\begin{frame}{Modelo Entidade-Relacionamento}
    \small
    \begin{center}
    \begin{tikzpicture}[
        entity/.style={draw, rectangle, rounded corners, fill=ClinicBlue!20,
                       minimum width=2.8cm, minimum height=0.7cm,
                       font=\small\bfseries, text centered},
        rel/.style={draw, diamond, fill=ClinicBlue!40, minimum width=1.8cm,
                    minimum height=0.6cm, font=\tiny, text centered},
        line/.style={thick}
    ]
        \node[entity] (users)  at (0, 2)    {users};
        \node[entity] (medicos) at (4, 2)   {medicos};
        \node[entity] (appts)  at (4, 0)    {appointments};
        \node[entity] (avail)  at (8, 2)    {doctor\_availability};

        \draw[line] (users.east)  -- node[above]{\tiny 1} node[below]{\tiny 1} (medicos.west);
        \draw[line] (users.south) -- ++(0,-1.2) -| node[near start, above]{\tiny 1} node[near end, left]{\tiny N} (appts.west);
        \draw[line] (medicos.south) -- node[right]{\tiny 1} node[left]{\tiny N} (appts.north);
        \draw[line] (medicos.east) -- node[above]{\tiny 1} node[below]{\tiny N} (avail.west);
    \end{tikzpicture}
    \end{center}

    \vspace{0.1cm}
    \begin{columns}
        \scriptsize
        \begin{column}{0.25\textwidth}
            \textbf{users}\\id, name, email,\\cpf, password\_hash,\\role, phone
        \end{column}
        \begin{column}{0.25\textwidth}
            \textbf{medicos}\\id, user\_id,\\specialty, crm, bio
        \end{column}
        \begin{column}{0.25\textwidth}
            \textbf{appointments}\\id, patient\_id,\\doctor\_id,\\scheduled\_at, status
        \end{column}
        \begin{column}{0.25\textwidth}
            \textbf{availability}\\id, doctor\_id,\\weekday,\\start/end\_time
        \end{column}
    \end{columns}
\end{frame}

% ============================================================
% SLIDE 11: REGRAS DE NEGÓCIO
% ============================================================
\begin{frame}{Regras de Negócio}
    \begin{columns}
        \begin{column}{0.5\textwidth}
            \begin{block}{Agendamento}
                \begin{itemize}
                    \item[\color{ClinicBlue}$\checkmark$] Um slot só aceita \textbf{uma consulta} por médico/horário
                    \item[\color{ClinicBlue}$\checkmark$] Slots são gerados dinamicamente com base na disponibilidade cadastrada
                    \item[\color{ClinicBlue}$\checkmark$] Apenas dias com horário disponível são exibidos ao paciente
                \end{itemize}
            \end{block}
        \end{column}
        \begin{column}{0.5\textwidth}
            \begin{block}{Ciclo de Vida da Consulta}
                \begin{center}
                \small
                \texttt{scheduled} $\rightarrow$ \texttt{completed}\\
                \texttt{scheduled} $\rightarrow$ \texttt{cancelled}
                \end{center}
                \begin{itemize}
                    \item[\color{ClinicBlue}$\checkmark$] Só consultas \texttt{scheduled} podem ser canceladas ou remarcadas
                    \item[\color{ClinicBlue}$\checkmark$] Remarcação = cancela atual + novo agendamento
                \end{itemize}
            \end{block}
        \end{column}
    \end{columns}
    \vspace{0.2cm}
    \begin{alertblock}{Integridade Referencial}
        \small
        Exclusão em cascata: remover um médico apaga automaticamente suas consultas e disponibilidades via \texttt{FOREIGN KEY ... ON DELETE CASCADE}.
    \end{alertblock}
\end{frame}

% ============================================================
% SLIDE 12: RESULTADOS OBTIDOS
% ============================================================
\begin{frame}{Resultados Obtidos}
    \begin{columns}
        \begin{column}{0.5\textwidth}
            \textbf{Funcionalidades entregues:}
            \begin{itemize}
                \item[\color{ClinicBlue}$\checkmark$] \textbf{4 perfis} de acesso distintos e isolados
                \item[\color{ClinicBlue}$\checkmark$] \textbf{6 médicos} pré-cadastrados com disponibilidade semanal
                \item[\color{ClinicBlue}$\checkmark$] Agendamento em \textbf{4 etapas} guiadas (médico $\to$ dia $\to$ hora $\to$ confirmar)
                \item[\color{ClinicBlue}$\checkmark$] Cancelamento e remarcação pelo paciente
                \item[\color{ClinicBlue}$\checkmark$] CRUD completo de médicos (admin)
                \item[\color{ClinicBlue}$\checkmark$] \textbf{15+ endpoints} REST documentados
            \end{itemize}
        \end{column}
        \begin{column}{0.5\textwidth}
            \begin{exampleblock}{Indicadores}
                \small
                \begin{itemize}
                    \item \textbf{4} tabelas no banco de dados
                    \item \textbf{7} arquivos HTML (páginas)
                    \item \textbf{6} módulos JavaScript
                    \item \textbf{1} API REST Express.js
                    \item Slots calculados para \textbf{14 dias} à frente
                    \item Interface \textbf{responsiva} sem frameworks externos
                \end{itemize}
            \end{exampleblock}
        \end{column}
    \end{columns}
\end{frame}

% ============================================================
% SLIDE 13: CONCLUSÃO
% ============================================================
\begin{frame}{Conclusão}
    \justifying
    O sistema \textbf{Clínica Asklepion} demonstrou que é possível construir uma aplicação web funcional, segura e organizada utilizando tecnologias acessíveis como \textit{Node.js}, \textit{Express} e \textit{SQLite}, sem dependência de frameworks de frontend pesados.

    \vspace{0.4cm}
    \begin{columns}
        \begin{column}{0.5\textwidth}
            \textbf{Pontos fortes:}
            \begin{itemize}
                \item Arquitetura clara em 3 camadas
                \item Controle de acesso por perfil
                \item Slots dinâmicos e sem conflito
                \item Fácil implantação local (XAMPP)
            \end{itemize}
        \end{column}
        \begin{column}{0.5\textwidth}
            \textbf{Próximos passos:}
            \begin{itemize}
                \item Criptografia de senhas (bcrypt)
                \item Notificações por e-mail/SMS
                \item Deploy em nuvem (VPS/Heroku)
                \item Prontuário eletrônico do paciente
            \end{itemize}
        \end{column}
    \end{columns}
\end{frame}

% ============================================================
% SLIDE 14: AGRADECIMENTOS
% ============================================================
\begin{frame}[plain]
    \centering
    \vspace{1.2cm}
    {\huge \color{ClinicBlue} \textbf{Obrigado!}}\\
    \vspace{0.8cm}
    \large Agradecemos a atenção de todos.\\
    \vspace{1cm}
    {\color{ClinicBlue}\Huge\textbf{+}}\\
    \vspace{0.4cm}
    \footnotesize \textbf{Clínica Asklepion} $\cdot$ Sistema Web de Gestão de Consultas $\cdot$ 2026
\end{frame}

\end{document}
